
\documentclass[../../main.tex]{subfiles}
\begin{document}

% 年度の大きな表示
\begin{center}
  {\fontsize{48pt}{58pt}\selectfont \textbf{1965}\normalsize\textbf{年度}}
\end{center}

\vspace{10mm}

% 本文

\vspace{15mm}

% 出題テーマと難易度の表
\noindent\textbf{出題テーマと難易度}

\vspace{5mm}

\begin{center}
  \shadowbox{%
    \begin{tabular}{|c|p{8cm}|c|}
      \hline
      \textbf{問題番号} & \textbf{テーマ} & \textbf{難易度} \\
      \hline
      第1問 & 図形と方程式（最短距離） & \\
      \hline
      第2問 & 三次方程式と解と係数の関係 & \\
      \hline
      第3問 & 三角関数の最大最小 & \\
      \hline
      第4問 & 積分公式（シンプソンの公式） & \\
      \hline
      第5問 & 積分と不等式 & \\
      \hline
      第6問 & 面積の最小値（微分） & \\
      \hline
    \end{tabular}%
  }
\end{center}

\vspace{15mm}

% 解答
\noindent\textbf{解答}

\vspace{5mm}

\begin{center}
  \shadowbox{%
    \renewcommand{\arraystretch}{1.6}
    \begin{tabular}{|c|c|p{8cm}|}
      \hline
      \textbf{問題番号} & \textbf{小問} & \textbf{答え} \\
      \hline
      第1問 & - & $P(3,2)$ \\
      \hline
      第2問 & - & $7a^2b$ \\
      \hline
      \multirow{2}{*}{第3問} & 最小 & $(x,y)=\left(\dfrac{\pi}{12},\dfrac{\pi}{12}\right),\left(\dfrac{5\pi}{12},\dfrac{5\pi}{12}\right)$ \\
      \cline{2-3}
                             & 最大 & $(x,y)=\left(\dfrac{\pi}{12},\dfrac{5\pi}{12}\right),\left(\dfrac{5\pi}{12},\dfrac{\pi}{12}\right)$ \\
      \hline
      第4問 & - & $\dfrac{1}{h}\int_{-h}^h f(x)dx=\dfrac{1}{3}\{f(h)+f(-h)+4f(0)\}$ \\
      \hline
      第5問 & - & 証明問題 \\
      \hline
      第6問 & - & $\dfrac{4}{5}\sqrt{\dfrac{6}{\sqrt{5}}}$ \\
      \hline
    \end{tabular}%
    \renewcommand{\arraystretch}{1.0}
  }
\end{center}

\vspace{15mm}

% 解答時間
\noindent\textbf{試験形態}

\vspace{5mm}

\begin{center}
  \shadowbox{%
    \begin{tabular}{p{8cm}}
      （要確認）
    \end{tabular}%
  }
\end{center}

\end{document}
